\documentclass[11pt]{article}
\usepackage[margin=1in]{geometry}
\usepackage[T1]{fontenc}
\usepackage{lmodern,microtype,amsmath,amssymb,array,booktabs,longtable}
\newcolumntype{P}[1]{>{\raggedright\arraybackslash}p{#1}}
\usepackage[hidelinks]{hyperref}
\hypersetup{pdftitle={Swarm Coordination Under Faults and Changing Conditions: Research Contract, Draft 0.3.1}}
\setlength{\parindent}{0pt}
\setlength{\parskip}{0.4em}
\setlength{\emergencystretch}{2em}
\newcommand{\repo}[2]{\href{https://github.com/eternal-november/eos-research-programmes/blob/4359a8a1484498b9c46d8e58a0112448a1fba06b/projects/#1}{#2}}
\newcommand{\Runs}{\operatorname{Runs}}
\title{Swarm Coordination Under Faults and Changing Conditions\\Research Contract}
\author{}
\date{Draft 0.3.1: coupled recovery and comparative evidence \quad 24 November 2026}
\begin{document}
\maketitle

\section{Governing question and research return}
How can groups of autonomous vehicles keep doing joint work while the conditions for coordination change? In particular, when a change in evidence, physical capability, authority, roles or goals affects an ongoing joint activity, what remains realizable from the participants' actual views and powers, what must survive the change, and what must be acquired, renegotiated, repaired or physically settled before useful work can continue?

The research return is a small body of preservation, replacement and impossibility laws that explain these changes across materially different coordination problems. The engineering return is one stack that preserves safety and useful service across stated interacting faults. Enclosed-space coordination is also a principal research target: derive when physically interdependent missions can continue through a change in both reachability and contact, then test that result on the integrated stack. A collection of independent handlers, a common notation, and a generic assembly theorem with the hard hypotheses supplied do not suffice. Conversely, the aim is not to force every control, causal, risk or strategic problem into one universal theorem. Retain distinctions and separate specialized results wherever their laws do not transport.

The programme studies joint requirements and the live commitments, effects and dependencies created in realizing them. Commitments include incurred duties and granted permissions whose effects or settlement remain relevant; they are not merely messages or mutable task labels. Joint viability may also constrain work before a commitment is incurred, and a law can concern a different object where that is more appropriate. The common problem is \emph{changing conditions of realizability}, not preservation of every original plan or allocation.

This is a new contract in the swarm-coordination project. The completed v0.5 contract is prior evidence and a source of regressions, not a list of obligations to reopen. The frontier stress tests are calibration: parity, fixed report schedules and set-cover instances already handled by ordinary methods do not count as research advances. Work is for civil missions and uses proofs, executable models and simulation; no physical tests or claims of deployment readiness. Earlier private vehicle research is not an input.

\paragraph{Preservation of obligations.} This revision preserves the substantive research and integration obligations of v0.1 while retaining and strengthening v0.2's shared-law organization. A common result may discharge several obligations where its instantiations establish the required extension, guarantee or resolving impossibility in each. Reduction to established results supplies inherited knowledge and may justify revising a research question, but does not by itself discharge a required extension. R1--R5 retain that extension burden; R6 retains its explicit exception for rigorous synthesis or comparison. Baseline reproductions, supporting lemmas and individual test cases need not each be novel. Domain-specific obligations remain wherever the shared laws do not suffice, and a result for one subclaim does not silently close the other subclaims of its track. Section~\ref{sec:audit} records the restorations, additions and retained scope. This preservation statement governs the present draft and its completion claims; it does not prescribe or freeze future amendment governance.

\section{Common semantic obligations}\label{sec:semantics}
Use a shared account of actual evolving histories, with explicit typed projections or compatible model interpretations. No particular state-machine formalism, policy language, centralized ledger or categorical representation is imposed. An abstract carrier is adequate only relative to the claims and consumers that use it. Establish the following distinctions before a proof, comparison or controller relies on them.

\begin{description}
\item[World and effects.] Physical state, trajectories, true event occurrence, observation, report, acknowledgement and realized effect are distinct. Retain effect identities, in-flight operations, physical settling and resource consumption where consequential. A failed, partial or interrupted operation has actual outcomes.
\item[Local information.] Specify the initial design knowledge, each actor's local history, admissible communication and memory, identity, provenance, freshness and timing. Recognition is not occurrence; a shared record is not necessarily common knowledge. Information acquisition can alter the world, future information and duties.
\item[Powers and owners.] Distinguish ownership, valid and recognized authorization, actual access, influence and enforcement. Identify the owner of each relevant duty, who can amend or discharge it, the actual retirement of old powers and continuing liabilities. Where strategic participants matter, state their preferences, participation constraints and enforcement separately from the joint mission law.
\item[Requirements and resources.] Identify the beneficiaries and holders of requirements. Distinguish hard invariants, actual required service, recovery deadlines, conditional duties, and legitimately amendable goals. Record shared and private reserves and dependencies, not independent copies of the same resource. A safety constraint, risk cap and authorization can have different amendment laws; a physical constraint cannot be amended into feasibility.
\item[Source alternatives.] Specify a nonempty joint class of initial states, fault/environment histories or causal strategies, and, where appropriate, fixed stochastic laws. Distinguish a faulty sensor, recipient, controller, actuator, enforcer and shared cause; harmful interaction can also arise under nominal local conduct. No independent per-role budget or true-fault oracle may be introduced by decomposition. Distinguish faults that destroy a requirement directly from failures of the response.
\item[History and time.] Preserve chronology, causal predecessor, local clock uncertainty, elapsed costs, pending and divergent work, and what survives reboot, rejoin, substitution or change of name. Record the horizon and workload, fairness and progress conditions. A timeout, deadlock or prematurely stopped maximal history receives a mission verdict; it is not discarded.
\end{description}

For a cooperative robust fragment, let $\theta$ specify the design, source family, law, admissible local policies and resources. A positive realizability judgment has the schematic form
\[
 \exists\pi\in\Pi^{\mathrm{local}}_\theta\;\forall s\in\mathcal S_\theta:\quad
 \left[\begin{gathered}
 \Runs_\theta(\pi,s)\ne\varnothing\;\land\\
 \forall h\in\Runs_\theta(\pi,s):\;\Phi_\theta(h)
 \end{gathered}\right].
\]
Here $\mathcal S_\theta$ is nonempty and independent of whether the controller succeeds, $\Runs$ includes admitted maximal executions (including failures and finite stops), and $\Phi$ retains safety, specified actual service and required lawful responses. Policy selection precedes the hidden scenario; causal environment reactions to the policy are allowed under the declared source law. The compliant actors implement $\pi$ from their own information and powers. A law governing their response does not make a faulty physical effect impossible by definition. Full-information feasibility, local realizability and realized execution are separate judgments. A per-history existence of a suitable action is not a locally implementable policy.

For stochastic obligations give the corresponding policy and probability quantifiers on actual trajectory laws, with conditioning and support stated; for competing owners give the strategy, compliance, enforcement and solution predicates. The displayed robust judgment is not a universal substitute for them. Each positive family must inhabit its primitive assumptions, including joint resources and fault combinations, without supplying the desired information, lawful repair or useful service as an oracle. If the source models differ across families, supply the actual interpretation and its guarantee direction; common names or matching endpoints are insufficient.

\section{Shared laws and the law-to-interface problem}\label{sec:laws}
The following three families of laws are the initial organizing targets. Their final mathematical statements, exactness, domains and number are to be discovered. They may split, merge or be reduced to existing results; the research obligation is to explain the cross-family structure and its limits, not to prove these formulations by stipulation. The shared research return requires substantial new results (constructive or impossibility) resolving declared questions across materially different research families. At least one substantive cross-family result is necessary, but is not sufficient to discharge the programme: the connected scientific return below and the substantive R1--R5 targets must also be met, with R6 retaining its synthesis exception. A single result may meet several targets through demonstrated transfer; no quota of separate theorems is imposed. The core engineering return also requires constructive, reused guarantees, which may be inherited. A negative result alone cannot authorize an architecture that fails the positive operating envelope.

\subsection*{L1. Realizable continuation and interface construction}
Given the current history, live duties, joint requirements, actual assembly and admissible changes to it, derive what the participants must be able to observe, do and agree to in order to realize a continuation. Determine when existing local views and powers suffice, and when a change in those interfaces is necessary. The output is a usable local policy or sufficient local implementation obligation with a source-valid construction, or a delimited obstruction with a boundary at which realizability changes. It is not an omniscient selection from feasible joint outcomes.

For a deterministic assembly, the inverse image of the joint requirement describes the acceptable joint histories. With nondeterminism or strategic or probabilistic choices, declare the appropriate quantifiers. A relation of acceptable joint histories does not itself factor into implementable local policies. Derive the additional observation, coordination, coupling, resource or authority constraints that make the assembly realizable, and establish whether the resulting choices can be made from local histories. Compare alternative architectures at the same joint law: e.g., acquiring evidence, deferring a choice, transferring an allocation or changing who can act. Charge the work and duties created by the intervention. Identify infeasibility arising from information separately from physical or normative infeasibility; do not infer non-realizability from the failure of one sufficient certificate.

The intended increment over existing local-policy synthesis and the Contextual Composition imports is construction or characterization of \emph{changing} executable interfaces for live work: earlier actions affect the later demand, possible evidence, and available authority or capacity. Candidate cross-family tests include R1's evidence cuts and R3's costly evidence, or R1 and R5's local authority obstruction. A different physical story over an unchanged decision table is not, by itself, a second structural use.

\subsection*{L2. Live replacement, residual sufficiency and dependency propagation}
Establish when a change to a mechanism, role, evidence supplier, authority or consumer preserves an outstanding obligation and when a receiver must instead revalidate, migrate, renegotiate or settle it. Index the guarantee by the admitted contexts, live histories, observations, policy class and costs. An unchanged summary licenses replacement only for the tests or continuations that actually descend through it. A directed refinement can suffice; do not assert equivalence, reverse reflection or arbitrary controller transport without evidence.

A useful diagnostic is contextual distinguishability. If two histories have the same retained account but an admitted consumer must distinguish them to establish a guarantee or choose a lawful response, a summary-only method cannot in general implement that consumer. The existing residual-profile construction gives a coarsest profile for a declared primitive testing interface, not automatically one for newly enabled consumers or arbitrary recurrent execution. The open burden is to derive the consumed tests and their native/source witnesses from the changing assembly, maintain adequate accounts under live changes, and construct migration or continued settlement paths when distinctions were not retained. Alternatively establish a sharp obstruction to such maintenance in a declared class. Keep reserve distinctions if their prospective consumers are not yet known.

The law must name the initial relation, the actual replacement operation and its outcomes, dependency discovery, the proof of preservation or invalidation, and the continuation or residual guarantee after the change. Establish compatibility under repeated or overlapping replacements within the admitted envelope. Separate semantic validity, efficient certificate validation and the cost of the entire repair. New obligations can be created by recovery; a proof of unchanged local outputs may not preserve a receiver's consumed historical dependencies. Candidate transfers pair R2 with a genuinely different supplier or consumer in R1, R4, R5 or R6; R2 retains its independent dependency-change and invalidation tests. Revalidate all hypotheses affected by the change; no receiver rewrite is allowed in a claim of unchanged-receiver proof reuse.

\subsection*{L3. Same-history joint admission, accounting and actual settlement}
Determine when local guarantees about obligations, resources, exposure and progress can be assembled into a joint guarantee on the \emph{same} source history. Track allocation, coupling, outstanding authorizations, incurred effects and the physical or digital events that actually settle them. Independent feasibility, marginal safety, valid signatures and a common abstract readiness flag do not in general compose. Use native comparison or direct common-history witnesses that retain the relevant dependencies, with the exact direction of guarantee transport stated. Derive the necessary joint admission and settlement premises from actual local mechanisms; merely postulating a sufficient joint envelope is not the constructive return.

Distinguish at least the following kinds of conclusion: admissible joint choices (including robust or stochastic constraints), exclusive or conserved actual effects where the task requires them, and progress to settlement within the stated service envelope. They may require different laws. Shared structure is claimed only where their source interpretations and composition obligations actually match. A sufficient risk or opportunity certificate is not an exact policy synthesis or a lower bound on all possible controllers. Conditional fairness is not an elapsed-time bound. A sound proof-side existential choice of a surviving route is not information available to the running vehicle.

Use the common-history, ownership and controlled-recovery results of Contextual Composition and the bounded exclusive-service rule of Byzantine Control where applicable. The intended new work is the maintenance or derivation of the needed joint interfaces as handoffs, evidence, allocations and physical recovery options change, including unsettled effects across a second fault. Candidate transfers pair R4's coupled exposures with a materially different actual-effect or resource family (R2, R5, R6 or I1). For risk and physical viability do not identify probability bounds with Boolean possibility sets, or capacity conservation with geometric feasibility. If their only common structure is preservation of a joint history, retain that limited result and develop the specialized obligations separately.

\subsection*{Cross-family and architectural acceptance}
The principal results must form a coherent body of guarantees and obstructions. For every research family, identify the consequential input, output, hypothesis, replacement condition or obstruction that it contributes to, receives from or establishes for that body and its consumers. Count a connection only when an actual interpretation transfers such a result and its omission would invalidate a claimed reuse or change an admitted continuation, architectural choice or required guarantee. Common notation, a shared state carrier, conjunction of separately proved claims and a common simulation do not establish this connection. Where stronger sharing fails, state the boundary, retain the specialized obligations, and establish compatibility between their supplied interfaces and their consumers. A pair of linked tracks with the remaining scientific work disconnected does not establish the programme-wide return. No single universal theorem or arbitrary graph-connectivity metric is required.

For each law offered as a substantive shared return, state a single quantified result with an explicit context and comparison class. Give at least two materially different instantiations with source-to-law maps, concrete establishment of the consumed hypotheses, and transport of the conclusion. For a new shared result, at least two research families must be served; an integration trace alone is not a second mathematical instance. State what remains specific to each family and what new dependency invalidates the law. The comparison is with both the closest existing results and an equally factored, cached direct account; a tie is a legitimate outcome. A shared data type, common lemma about conjunction, or relabelled existing theorem is not a substantive new cross-family result.

At least one constructive shared guarantee must constrain actual recovery or continuation decisions in the integrated stack, and at least one new shared result must materially affect such a decision or an architectural choice that enables it. The same law may meet both requirements. Show the consequential difference on matched histories. A new impossibility law can instead explain a refused operation, a necessary interface change or an exclusion boundary with a corresponding positive design; it does not become a runtime oracle. The architecture need not implement a novel formal calculus if a simpler verified mechanism supplies the same guarantee. The technical return should explain why the adopted mechanisms are sufficient, when they are not, and what would have to change.

\section{Existing research and reduction discipline}\label{sec:imports}
The two related project programmes are inputs at their actual theorem and source scope. Their results are not fresh swarm claims and their open research directions are not completed imports. Section~\ref{sec:sources} supplies precise retrieval routes at the repository revision used for this draft. Read the relevant statements, proof or checked declarations and counterexamples before assigning an inherited guarantee.

\paragraph{Contextual Composition.} This is the named relative construction of independently specified context and content with dependent interaction and native comparison maps, not a generic synonym for compositionality. The \repo{contextual-composition/publications/interactads/Interactads.tex}{Interactads manuscript} supplies a mixed calculus and a gluing theorem under actual lifting, pasting and pure-interface hypotheses. The \repo{contextual-composition/feedback-and-concurrency/observable-covers/Residual_Profiles_and_Admission_Reflection.tex}{residual-profile criterion} constructs adequate profiles for declared primitive tests, with reflection and attained-class coverage required for a common presentation along fixed native witness maps. It does not construct arbitrary new consumers, a small recurrent summary or a coherent policy selector. The \repo{contextual-composition/protocols/refinement-recovery/Contextual_Refinement_and_Controlled_Recovery.tex}{controlled-recovery and refinement theorem} supplies sufficient finite contextual comparison and residual guarantees; its scope is not unrestricted concurrent physical recovery.

The \repo{contextual-composition/foundations/native-whole-scope-feedback/Nonblocking_Recurrent_Assembly_Reading_Route.tex}{nonblocking recurrent assembly route} distinguishes global admission, calibrated whole-law control, and decentralized realization. Its retained commitments and historical effects must not be replaced by fresh independent residual draws. Its mixed comparison and progress transport use actual source/control hypotheses, not equality of endpoints. Import those that fit; if the desired interaction changes the native grade, observer, authority or available continuation, supply the missing lift, comparison or source semantics. The \repo{contextual-composition/research-directions/from-joint-requirements-to-interfaces.md}{law-to-interface research direction} explicitly leaves construction and native transport of such interfaces open. It is a close research antecedent to L1--L2, not evidence they are already solved. Use \repo{contextual-composition/research-directions/designing-concurrent-interactions.md}{the concurrent-activity design route} to compare alternative arrangements and their residual duties.

\paragraph{Byzantine Control.} The \repo{byzantine-control/portability/COMPLETION-AUDIT.md}{reusable control results} already report a common-history recurring useful-return proof across two different positive architectures, recipient-relative retirement, source-local adaptation and a held-out change. They distinguish observation, authorized action and actual independent effect; a release can leave live powers, work and liabilities behind. They also report a comparison where an equally factored direct account discharged the same global steps. Do not claim novelty for these ingredients or presume that framing alone reduces proof effort.

The \repo{byzantine-control/exclusive-service/RULE.md}{bounded exclusive-service rule} gives one actual useful effect by a deadline on every admitted maximal history when the whole-window denial envelopes are covered by locally justified routes, each committed issue settles in bounded time, and source-specific exclusion holds. It is a conditional certificate, not a general synthesis theorem; its finite cover is ordinary availability/transversal mathematics and a failed route-library cover is not a universal impossibility. Attempt this reduction for swarm report, custody, clearance or recovery effects where exclusivity is part of the law. The \repo{byzantine-control/DOWNSTREAM-CONTRACT.md}{downstream contract} and \repo{byzantine-control/REMAINING-PROGRAM.md}{remaining programme} distinguish retained results from further research duties. Reuse local information, actual powers, joint fault support, lawful intervention, recurrent history and renewal results only where the actual premises transfer. Multi-owner cooperation and strategic conflict are not synonymous with Byzantine failure; the swarm contract does not require reproducing the entire Byzantine programme.

\section{Acceptance and evidence}\label{sec:acceptance}
For every research obligation and each materially distinct subclaim, identify the nearest result at theorem, algorithm or implemented-mechanism level. State its hypotheses, map it to a shared instance, and attempt the strongest known reduction before freezing the extension claim. Record what is inherited and what changes: an assumption, guarantee, cost or class of instances. A citation, new notation or larger simulation does not establish progress. If an existing result already covers a research target, retain it as a baseline and revise the research question openly; distinguish the inherited resolution of the earlier target from the remaining or revised extension claim. An adequate engineering mechanism need not be made artificially novel, but its adoption does not convert an unmet research extension into a completed obligation. The shared substantive return in Section~\ref{sec:laws} supplements the R1--R5 research burdens; it does not replace them. R6 retains its stated synthesis exception.

For each evaluated claim record the precise proposition or acceptance predicate, family and exclusions, source and policy quantifiers, stakeholders and amendment powers, admitted faults and resources, concrete suppliers of evidence, baseline and exact proposed increment, cost and evidence instrument, counterexample or falsifier, and how the conclusion is consumed. A family needs a nontriviality witness showing that the proposed differentiating phenomenon actually changes a guarantee, feasible policy, admissibility decision or cost. Review the family boundary itself, not just the cases selected from it. Each claim must distinguish inherited resolution, new result, rigorous synthesis, scoped impossibility, partial result and open work. A failed technique, failed sufficient condition or scoped lower bound is partial unless it resolves the declared claim. An impossibility must name the class, violated requirement and constructive boundary where feasible. It closes a target only when it resolves that target at its declared scope; the constructive-or-impossibility alternatives in L1--L3 do not replace a separately required construction, experiment, positive witness or service guarantee in R1--R6 or I1--I2. Stopping work on an open question is not its scientific resolution.

Freeze each evaluated family used for confirmation or closure, fault envelope, comparators, resource accounting and acceptance predicates before independent case selection. Exploratory cases may inform a revision, but may not be relabelled as independently selected confirmation; a case-specific change requires fresh independent cases for the changed claim. Exploration and revision of the contract remain open; record the actual version and scope of each result and do not present evidence from a changed problem as a result for its predecessor. Declare which observations arrive to which actor before each irreversible choice; common knowledge, synchronized clocks, reliable replacement and an omniscient planner are not defaults. Compare at matched information, authority, physical resources, workload, deadlines and policy expressiveness. Charge offline preprocessing, retained state, dependency validation, online computation, communication and physical opportunity costs where they support a claimed gain. A full-information solution may diagnose a gap, but cannot serve as the implementable comparison. Report safety and useful service separately; idling, blanket refusal or silently draining the mission cannot discharge a service requirement.

\paragraph{Advancement and comparison purposes.} For I1 and I2 distinguish (i) operational advancement, where a constructive guarantee and source-valid matched failure or exclusion establish capability unavailable to the stated reference, (ii) assurance advancement, where stronger guarantees can be established without showing a difference in actual capability, and (iii) efficiency or service advancement at a matched mission and operating envelope. If $E_A$ is the actual envelope of a fixed architecture under the declared law and $\widehat E_A$ its soundly established envelope, then $\widehat E_A\subseteq E_A$. A strict enlargement of established envelopes alone does not establish a strict enlargement of actual envelopes. An admitted history that defeats a fixed reference does not rule out all implementations of its method; a class-wide claim needs a class-wide argument. Compare bounds with the same meaning, quantifiers and evidence scope: a tighter certified upper bound is not necessarily a lower actual worst-case cost. Assurance-only improvement is valuable but does not by itself discharge the positive architectural or coordination advancement required in I1 and I2. A material operational, resource or service advance must preserve non-waivable requirements and respect all resource constraints; identify and justify any worsening elsewhere against the declared stakeholders' priorities. Do not label a trade-off as dominance or use an undeclared aggregate objective.

Keep three comparison roles distinct. The strongest applicable established mechanisms, including their reasonable combination, test the substantive research and capability increment. An equally factored, cached direct account or implementation of the proposed mechanism tests the additional value of its contextual formulation; it may tie or win without negating a genuine advance over established mechanisms. A dependency ablation tests necessity of a retained distinction or supplier and is not an equal-interface performance comparator. Match primitive information, powers, resources, workload and deadlines; charge any additional evidence, access or capacity constructed by a candidate, and state which changes to the interface are the claimed advance. Establish the relevant guarantees and source correspondences in each arm. Report a failed reduction or an unavailable implementation at its actual scope, not as a universal limitation of existing methods.

Use Contextual Composition substantively on at least one named interface and a named preservation or replacement law. Identify the duties and observations allowed to cross that interface, its consumers, native objects, comparison strength and retained witnesses; establish the law's hypotheses from the concrete system and derive a reused result. Distinguish evidence retained for the mathematical argument from information that each running participant may actually receive, disclose or use. Additional construction or obstruction applications do not replace this preservation-or-replacement obligation. An application result requires actual interpretation, not reuse of terminology or moving a proof into a wrapper. Compare with an equally factored, cached direct method and report ties or superiority. The architecture may use the better implementation; substantive mathematical use of Contextual Composition does not require a runtime implementation of a categorical calculus. No universal native gluing, general decentralized selector, proof-effort saving or categorical novelty is assumed. Retain source versions, matched reference implementations, counterexamples and reproduction commands. Use Lean for stable, load-bearing proof kernels where it adds assurance; keep the abstraction and source correspondence separate from the checked kernel and from simulation evidence.

\section{Research families: discriminating tests of the laws}\label{sec:families}
R1--R6 specify families and discriminators, not six independent architectures or six predetermined theorem statements. Each family must identify which shared law it consumes, supplies, extends or obstructs, and which additional domain-specific results it needs. A result may serve several families only through an actual interpretation and proof of the transfer obligations in Section~\ref{sec:laws}. Do not repeat an inherited theorem merely to give each family its own result. Consolidation changes the organization of the proofs, not the substantive burden of the targets. The following targets retain their own acceptance conditions, including separate subclaims, constructive outputs, comparisons and held-out tests.

\subsection*{R1. Coordination across communication and observation regimes}
Determine when a change in evidence forces a change in the joint maneuver (L1), and which retained distinctions let a receiver preserve or change its response (L2). Start with the project communication reference and \href{https://xiangyin.sjtu.edu.cn/Paper/15Automatica.pdf}{dynamic coobservability}. Give a rule, with a sharp counterfamily, for preserving or replacing a live commitment across two different evidence cuts that lead to the same degraded regime. Include one-way, delayed, partitioned and absent contact, including no usable observation or prior exchange for the decision. Test channel substitution with a held-out non-radio modality under the same task requirements. Treat passive vision separately: establish what anonymous or intermittent tracks suffice for, and when named duties require persistent identity or a physical interlock. Preserve distinctions among message reception, action evidence, and mechanical exclusion. Declare what preinstalled conventions, environmental observations and physical interlocks remain usable even when no messages arrive. Separate a necessary intersection of locally admissible actions from the existence of one compatible family of continuation policies.

A pair of histories indistinguishable to the deciding actor but requiring incompatible actions is a genuine obstruction only for the declared information, authority, deadline and policy class. Identify a feasible change to that interface where possible; compare the change against dynamic-observation synthesis at the same interface.

\subsection*{R2. Repeated repair, propagation and live synthesis}
This is the primary live-replacement test of L2, consuming L1's continuation judgment and L3's unsettled-effect account. Start from \href{https://pmc.ncbi.nlm.nih.gov/articles/PMC9468080/}{LiveLTL obligation monitors} and the \href{https://smlab.cs.tau.ac.il/syntech/dupdate/dupdate-icse22-cr.pdf}{dynamic update synthesis construction}. Extend the baseline to a family with two overlapping handoffs, delayed recognition, and live duties, where the first lawful repair changes which later crew can acquire decisive evidence or authority. Include role and partner changes, split and rejoin without common reset, and a case in which a second local repair either propagates or is contained depending on a retained dependency.

Prove which certificates survive a replacement and which must be recomputed; establish a bound on affected work and retained state under a declared recurrence envelope, with a family witnessing an essential dependency. Freeze one receiver proof and check a held-out change that actually changes a consumed dependency, with no receiver rewrite. Compare to a factored cached observer and strong whole-system reference. Evaluate preinstalled policy and bounded online synthesis at matched resources, including any service forgone while it runs. The safety path must remain valid if computation is late or fails; characterize the service actually preserved, including work by unaffected crews while repair is pending. A global lock is not a free safety mechanism.

Include dependency discovery, migration, invalidation checks and maintenance in the affected-work bound; a graph traversal over supplied dependencies alone does not establish the claimed bound. Compare the reuse of a valid certificate, the cost of verifying continued validity and the cost of the complete replacement separately. Reuse the existing residual-profile and controlled-recovery results only under their respective hypotheses; the new burden is to obtain and maintain those hypotheses during the handoffs. No speedup is required over an unrestricted exact solver, but a reuse or cost claim must survive the matched comparison.

\subsection*{R3. Actionable evidence and anticipatory control}
\paragraph{R3a. Creating actionable evidence (L1, L3).} For costly evidence, build from \href{https://arxiv.org/abs/1209.1695}{common-information dynamic programming} and \href{https://arxiv.org/abs/1410.0083}{reactive synthesis with active sensing}. Characterize a class where the current choice changes which observations another actor can lawfully obtain before its next decision. Charge the action and resulting duties; compare with passive control and the strongest matched active baseline. Exhibit a case where a more accurate but later or privately held observation cannot support the joint intervention. Show that the required observation is acquired from primitive, fallible operations rather than supplied as an already useful inference; retain the information and physical effects of a failed attempt. Where an actor can select a joint law or intervention, distinguish actual readable data, authority to select, and availability of a compatible execution mechanism.

\paragraph{R3b. Identification and preparation (L1).} For forecasts, test a separate question against \href{https://proceedings.mlr.press/v286/marmarelis25a.html}{policy evaluation with hidden trajectory confounding} and \href{https://arxiv.org/abs/1205.0858}{controlled sensing}. Show what is identified from selectively observed histories when earlier commitments affect which later experiments are feasible. Require a non-identifiable pair of models and either a sharp bound or a constructive intervention that separates the relevant future decisions at a stated cost. Reproduce the fixed-experiment case; independent repetitions of a fixed pair test or a forecast that only triggers feedback at the time of failure do not discharge this target. Declare the estimand, candidate source laws, available experiment policies, observation distributions, and downstream decision loss. The non-identifiable pair must agree on what the allowed data collection reveals and disagree on a decision-relevant feasible counterfactual; distinguish bounds relative to the existing experiment from the effect of designing a new experiment. Report what anticipation changes before the fault and the cost of a false forecast. R3a and R3b have separate claims and evidence; successful active sensing alone does not close causal identification.

\subsection*{R4. Joint risk and changing goals}
\paragraph{R4a. Joint admission (L3).} Build from \href{https://groups.csail.mit.edu/mers/papers/Ono-Williams_AAMAS10.pdf}{distributed risk allocation} and \href{https://arxiv.org/abs/1811.00215}{robust dynamic control with coupled uncertainty}. For a scalable family with overlapping exposures, derive an exact local admission criterion or a sharp limit on how much must be shared to implement one. A single fixed but uncertain law governs the whole mission; maintain the original mission-wide bound on the probability of any specified failure. Do not quietly apply a fresh cap to each decision or replace the uncertainty by independently chosen worst cases. Compare against factored inference and a common-information controller; marginal-only rules are negative controls, not the strongest baseline. Declare the event, horizon, policy-induced trajectory law, ambiguity set, and conditioning/update rule. An ex ante bound and a bound conditional on every reachable observation are different specifications. Distinguish exact admission, conservative admission and communication or storage lower bounds. A fixed law with latent dependence is not an instruction to resample a residual law after repair. Where applicable, attempt reduction to the calibrated whole-law and retained-dependence results of Contextual Composition.

\paragraph{R4b. Revision with live exposures (L2, L3).} Separately specify what happens when an authorized goal or risk limit changes while work is unresolved. Test both a legitimate prospective exemption and a change that binds live exposures, under timely and late notice. Account for commitments that cannot be withdrawn and show when reallocation preserves service and when it cannot. Distinguish an amendment to future requirements from a claim about satisfaction of the original mission law. A valid prospective exemption does not erase an incurred physical exposure; a retrospective change must not be reported as preservation of the original guarantee. A new cap or goal cannot be chosen after observing which policy passed.

\subsection*{R5. Changing authority, several owners and conflict}
Use L1 to derive locally realizable powers and L2--L3 to account for their changes and actual effects. Begin with the applicable results of Byzantine Control (Section~\ref{sec:imports}), not with a fresh static consensus problem. Use also \href{https://doi.org/10.1007/BFb0055622}{alternating refinement} and \href{https://pages.nyu.edu/debraj/Papers/KonishiRay.pdf}{dynamic coalition formation}; select the closest executable governance or game model for the frozen subclaim. Derive when an authorized amendment preserves existing rights and how the amendment changes who can authorize the next one. Treat four regimes distinctly: cooperative fleets without a common controller; limited trust among compliant owners even when reports are truthful and mutually consistent; inconsistent evidence or bounded faulty components; and actors with genuinely competing objectives. For the limited-trust regime, specify which reliance, disclosure, verification and refusal are permitted, and show when these limits change a realizable outcome or required interface. Distrust must not be represented solely as a faulty-component count, inconsistent reports or competing objectives. A common trusted controller or arbiter cannot be assumed for the decentralized cooperative regime. Select an administrative authorization model, including administrative role-based access control where applicable, for amendments of powers. For the conflict branch, declare the parties' preferences, feasible commitments, information, enforcement and solution concept; existence of a lawful joint plan does not prove voluntary participation, stability or execution. An agreement does not waive a non-waivable safety duty. Specify who can refuse, inspect, override or settle a duty, including human takeover where applicable. A minimal-cost selection of fixed consent sets is a baseline. Require a case in which changing the future powers changes a realizable outcome, and an independent obstruction when global agreement is possible in principle but local lawful execution is not.

Distinguish permission, recognized permission, physical access, actual influence, observation integrity and independent effects. Reuse Byzantine Control's source-sensitive support and retirement arguments and bounded exclusive-service theorem where their premises hold; the new multi-owner burden is changing legitimate powers, lawful refusals and jointly feasible settlement under changing evidence. A received takeover instruction, its authorization and effective transfer are distinct. A stale old command or an in-flight effect can survive revocation; a prohibition is not a physical gate. If an evolving fault or contamination model is consumed, preserve its history and derive restoration, rather than assuming clean re-entry. Do not import the entire Byzantine Control research programme as an additional swarm obligation. Keep abstract fault and authority models separate from intrusion or adversarial deployment.

\subsection*{R6. Fault-tolerant combinations of formation invariants}
Use L1 and L3 to distinguish lawful and informed recovery from joint physical viability; test L2 where geometry or capacity changes invalidate a receiver's guarantee. Extend or reproduce \href{https://doi.org/10.1016/j.jfranklin.2023.05.013}{Maeda, Endo and Matsuno's crash recovery}; contrast it with \href{https://arxiv.org/html/2609.29939v1\#S3.SS2.SSS3}{managed make-before-break departure}. Where control limits make individually feasible invariants incompatible, compare with \href{https://arxiv.org/abs/2210.01354}{composition of multiple barrier functions}. Study geometry, density, coverage, connectivity and clearance as mission requirements rather than a fixed shape. Use a family where at least two requirements compete for the same maneuver or spare resource after a fault. Distinguish invariants that must hold throughout, properties that can be restored by a deadline, and objectives that may be relaxed only through explicit authority. Establish a joint feasibility boundary or a verified recovery guarantee with a matched infeasible case, distinguishing physical infeasibility from inability to choose a recovery on local evidence.

State the joint viability and deadline-recovery conditions, disturbances, actuation constraints, time discretization and the soundness of any finite abstraction used by the coordination layer. An instantaneous feasible input does not prove a viable trajectory; distinguish a fault that directly destroys an invariant from a controller's avoidable violation and from the containment possible afterward. Restoring the old shape is not mandatory. This track may close as a rigorous synthesis or comparison without a novelty claim; it may not count suspended safety or silently abandoned work as recovery.

\section{Integration obligations}\label{sec:integration}
The architecture is the operational test of the shared laws, not a second programme of bespoke demonstrations. Each consumed theorem must have an implementation-level interface: the observations and resources supplying its hypotheses, the actual admission, repair or settlement decision it constrains, and the evidence that connects its output to the physical or mission-level guarantee. A proof-only witness is not a runtime input. The architecture may use a simpler established mechanism in place of a proposed new law if it provides the same required guarantee; record the disposition.

\subsection*{I1. Inventory-wide fault-tolerance stack}
Construct a single executable reference architecture and an assurance case for both project fault inventories. Pin inventory revisions, retain original entry identifiers, and map mechanism, cleanroom, extension and assurance entries without erasing mismatches. For each operational entry give a supported recovery or containment claim, its envelope and residual duties, and an executable representative; or give a validated exclusion or a falsifiable impossibility witness and what service is forgone. Unmodeled or untested entries remain open. Architecture-dependent entries require an applicability argument. Assurance entries, including wrong requirements and ineffective supervision, require design or audit checks with evidence. Coverage of entries means accounting for each, not universal tolerance or a promise to withstand arbitrary combinations. Distinguish tolerated, recovered within a bound, contained with service loss, inapplicable, impossible in the declared model, and open.

Before confirmatory evaluation, declare a positive integrated operating envelope independently of the inventory dispositions: mandatory service and safety, admitted primitive and compound faults, resources, workload and deadlines. It must include productive operation under communication loss, at least one repair that preserves a live duty, and a feasible compound-fault recovery with a second fault before settlement. Exclusions or correct refusals outside that envelope do not count toward its positive service floor. The numeric envelope and implementation are research selections, not fixed by this draft.

Build off the actual mechanisms in \href{https://arxiv.org/html/2103.11470}{NeBula} and, where applicable, \href{https://doi.org/10.1109/ICRA55743.2025.11128674}{HeRo}; establish full-text scope before making claims about HeRo. Start from their actual operating and recovery interfaces. State separately what is detected, what evidence authorizes a response, who owns outstanding work, what has been physically settled, and what the architecture may conclude about resumed service. Specify when a layer may rely on a recovery exported by another layer. Separate detection, authorized intervention, actual execution, residual-duty reconciliation and resumption as logical obligations, not necessarily serial runtime stages. No unqualified recovery flag or agreement alone establishes physical settlement. Use a common source history for each cross-layer proof: the local views, fault support, time, actual effects and resource use cannot be assembled from independently successful runs. Define the simulator-to-model correspondence and explicitly distinguish proof-side global state from controller-visible state. A second fault during repair, a common cause, and a fault masking another must be exercised with shared controller code. In particular, run matched restarts with and without an unresolved old authorization behind the same visible symptom. Independent selection of an additional compound case follows the interface freeze. Demonstrate a reused cross-layer safety and service argument on two materially different missions; compare to a typed, factored direct architecture at the same resources.

\paragraph{Required architectural return.} Give a reusable rule for composing recovery actions when one action changes the evidence, resources, or authority another action relies on. Derive its admission and resumption checks from the mechanisms serving those obligations. The rule shall preserve the named safety invariant and a positive service floor for every admitted maximal history in a declared compound-fault family, including an overlapping recovery in which the later fault changes the mechanism or participant that can settle an earlier duty. Prove this at the reference-model level, with a constructive connection to the executable stack. A successful replay of selected incidents is insufficient. This requirement does not extend to every combination of inventory entries or require an unbounded-history guarantee.

Use a common checking and repair mechanism on two materially different missions. At least one transfer must change the recovery dependency or the type of resource through which it is realized; relabelling the same transition system is insufficient. Freeze the mechanism before selecting the transfer. Named configuration parameters and adapters are allowed; a new recovery rule specialized to the held-out case requires another independent transfer. Exercise a case where each local handler can satisfy its own advertised result but the handlers' composition violates the joint law; show exactly which check blocks or repairs the composition, and include a positive case in which that check permits useful service.

For the frozen transfer, identify the unchanged decision procedure, reusable proof and guarantee, as well as the exact source-specific interpretation, adapter operations and proof obligations. Account for dependency discovery, compatibility validation, retained state, computation, communication and physical work in the adapter as in the common mechanism. Distinguish reused proof work from new source-specific arguments; no measured reduction in researcher effort is inferred merely from a frozen proof. An adapter supplying an already solved recovery, an unearned observation or a successful joint admission as an input does not establish reuse of the difficult decision. Show that the unchanged mechanism still derives a consequential admission, repair or resumption decision from the new source's actual evidence and powers.

For the local-handler counterexample, state the advertised individual guarantees, their hypotheses, context and quantifiers, and why each holds at its original scope. Identify the violated joint obligation and whether the failure is a premise invalidated by interference, a joint requirement not entailed by the local postconditions, an omitted consumed dependency, or a change outside the certified context. Derive the appropriate repair and establish the premises on the actual joint history. An intentionally weakened or ill-specified handler is a diagnostic control, not sufficient evidence of advancement over the strongest applicable established assembly.

Compare with the strongest applicable assembly of published recovery and planning mechanisms, together with a factored direct implementation preserving the same distinctions. Positive architectural advancement requires a material improvement in the admitted fault-and-workload envelope, or a strictly better declared resource requirement or service bound at a matched envelope. Attribute the improvement to an architectural decision by a matched comparison. A baseline may win, tie, or be adopted; if it absorbs the proposed advance, preserve the result and leave this advancement obligation open until a stronger target is met. Apply the advancement and comparator distinctions in Section~\ref{sec:acceptance}: a direct implementation of a genuinely new mechanism may supply the positive architecture even if the contextual implementation has no additional performance advantage; a tie with established mechanisms at the claimed advance leaves that advance unestablished.

For the shared constructive laws actually used, trace decisions on a common execution to the concrete suppliers of their premises. Use matched negative controls (e.g., omit the retained old authorization, collapse the coupled exposure, or assume a returned report is a settled effect) that change a consequential admission, repair or continuation. Include an independently selected receiver-compatible change and a legitimate invalidation that the stack detects. Compare against direct mechanisms with equivalent retained distinctions; ablations of required information diagnose necessity, not a fair performance baseline. Coverage tables, shared data structures and separate successful handlers do not discharge integration.

\subsection*{I2. Enclosed-space coordination and mission on the frozen stack}
Deploy that same reference stack in simulation for a civil inspection mission. Mission requirements include completed inspections, delivery of designated reports by a deadline, and verified egress or a defined fallback for mobile units under the admitted fault envelope. Choose report custody and vehicle recovery separately. Use a known interior layout to isolate partition effects, then a branching layout whose relevant map is learned in mission. Track physical routes, contact, line of sight and each local map independently. Hidden fault identity, true map, global clocks or custody status may be used by the evaluator but not supplied to a controller unless its interface actually makes them available. Build the benchmark from \href{https://superintelligence.jpl.nasa.gov/public/documents/papers/PlanRob-2020-Planning-and-Optimization-for-Multi-Robot-Planetary-Cave-Exploration.pdf}{JPL's cave planner}, with matched joint return-and-rendezvous as the comparison; \href{https://arxiv.org/abs/2605.20917}{SubTGraph} is a source of additional simulated layouts.

\paragraph{Required coordination return.} Study a family of enclosed inspection problems in which recoverable changes to physical reachability and communication jointly change which agents can finish time-critical work already underway. Include a topology change that redirects a feasible continuation rather than only delaying it; include a link partition during which useful work must continue and at least one unfinished obligation still couples agents on different sides. Show that the benchmark cannot reduce to independent exploration followed by eventual data upload, by exhibiting an admissible history on which independent plans fail the joint service law. State which changes affect sensing, motion, custody or settlement; geometry alone does not determine radio or visual contact.

The compound discriminator must be consequentially coupled: the reachability change and the contact change must both affect the same outstanding continuation obligation or a common admission, repair or settlement decision, not merely occur during the same mission. Trace their effects through actual routes, available observations, authority, resource and time constraints to that obligation. Compare the nominal, reachability-only, contact-only and joint-change cases wherever these are admitted at matched primitive conditions. If a combination is physically impossible or incompatible with the declared source law, identify the obstruction and use source-valid comparisons instead. Establish what remains to be resolved when the separately justified mechanisms are assembled; neither statistical superadditivity nor independent variation of physically coupled faults is required. Failure of an independent exploration policy is the required negative control; a stronger impossibility claim about independent policies must declare its class and establish the obstruction for that class.

Establish a constructive result for a declared family, rather than a controller tuned to one map. Derive a sufficient locally implementable continuation rule and a sharp or otherwise discriminating boundary: exhibit admitted histories that require different choices or feasibility verdicts, and establish what change in evidence, access, reserve or time closes that gap. An oracle may diagnose a boundary but cannot supply information to the working controller. Allow useful fallback in feasible cases; do not let blanket refusal, waiting for global reconnection, or reallocating all work before the fault discharge the target. Require a material advance at matched information, authority and resources over the best applicable coordination baseline, rather than merely a stronger application story. If the first family reduces, retain it as calibration and select a stronger family. An obstruction without a positive construction may close a question but does not discharge I2.

\paragraph{Baselines and attribution.} In addition to the existing cave planner, evaluate applicability of \href{https://arxiv.org/html/2103.11470}{NeBula's} return-to-contact and relay or data-mule mechanisms, \href{https://doi.org/10.3390/app16031212}{the mining workflow's} coordination under network interruption, and \href{https://whoenig.github.io/publications/2019_RA-L_Hoenig.pdf}{action-dependency-graph execution} for unexpected route disturbances. Also assess, as reduction candidates when the selected family consumes the corresponding interfaces, \href{https://doi.org/10.1109/TRO.2023.3344051}{switchable action dependency graphs} for dynamic reordering, \href{https://pmc.ncbi.nlm.nih.gov/articles/PMC10243298/}{epistemic planning under intermittent contact} for cooperative tasks and \href{https://doi.org/10.1109/TASE.2025.3558977}{temporal-logic allocation with inter-task constraints and variable capabilities}. These are comparison leads, not claims that their guarantees transport to the compound family. Use and combine the strongest applicable mechanisms, or demonstrate why they cannot express the matched case without changing their information or execution assumptions. Compare to the same combination augmented with ordinary factored recovery where applicable. Show which guarantees or gains depend on geometry-aware coordination, which on the stack's recovery rule, and which on their interaction. Where admissible, compare established coordination with established recovery, new coordination alone, new recovery alone, and their joint use under the same source law, mission and resource accounting. If a mixed variant cannot be realized without changing an interface or losing a premise, establish that incompatibility rather than inventing an executable comparator. Do not assume additive attribution. Trace the consequential dependency, for example when a recovery makes a new route executable and that route makes settlement possible. Apply the advancement and comparison distinctions in Section~\ref{sec:acceptance}. A shared result can satisfy I1 and I2 if it establishes each separately.

Freeze stack version, acceptance thresholds, permitted configuration, and comparators before independent choice of held-out geometry and event order. The geometry must change reachable recovery options, not merely visuals. Require an admitted episode combining blocked travel and lost contact with a second fault before settlement, a paired case where useful work continues through lost contact, and a boundary case where the stack gives a correct refusal or constrained fallback. Measure completed mandatory work, report deadlines, egress or authorized fallback, and reserves spent, not merely packet delivery or hazard flags. Compare to no-fault performance and the matched planner. A claim of improvement needs a declared outcome or resource advantage without trading away required safety. If the discovery or fault model admits unrecoverable trapping, explain the boundary rather than pretending that a fallback can always be executed. Case-specific code added after unblinding requires another holdout. State the evidentiary scope: a successful episode proves existence, not universal mission reliability; finite exhaustive checking or a proof applies only to its admitted model and a justified abstraction/implementation relation. If reporting empirical rates, declare the case-selection or sampling law and uncertainty. I2 does not replace I1 coverage or the proof of the shared laws.

\section{Moonshot: local-to-global synthesis for repeated handoffs}\label{sec:moonshot}
The moonshot asks for an effective, sound and complete compositional synthesis decision procedure on a restricted but nontrivial family of corridor layouts with $n$ segments for every positive integer $n$. Given a finite description of an admitted instance, the procedure must terminate and assemble the declared boundary summaries into implementable decentralized controllers, or correctly decide that no admissible controllers exist in the declared policy class. Establish the soundness and completeness of this summary-based assembly, including construction of the local controllers; bounded summaries alongside an unrelated whole-system synthesis procedure do not discharge the target. Its interface and composition laws must be uniform in $n$, with boundary summaries bounded by local parameters independently of corridor length and number of completed handoffs. This is uniform instance synthesis, not the stronger claim that one controller template works for every $n$, and not a guarantee that every admitted instance is realizable. A parameterized template theorem, if investigated, is a separate claim.

The family must permit overlapping neighboring handoffs, a bounded number of live commitments per boundary, and a lawful handoff that changes who can establish the next clearance. Service covers a prescribed class of offered jobs; refusing all jobs is not a feasible policy. State the finite instance representation, controller expressiveness (e.g., finite-state or a more general specified class), information model, policy memory, fault model, timing and local physical and evidence bounds. Justify the admissible policy class as part of the task family independently of the solver's success. Completeness for a restricted memory bound, controller grammar or other subclass does not establish non-realizability for more expressive controllers that remain admissible for the task. An algorithmic search bound alone is not a task restriction; any such narrowing changes the claim and must be reported as a restricted result rather than silently discharging the original target. Do not bound mission duration, impose a global reset or assume synchronized global observation to obtain the result. Declare sustainable offered workload and fairness or progress assumptions independently of a successful controller, and account for finite maximal executions. State how local memory, communication, and the end-to-end service deadline scale. Prove that replacement by a component with the same declared boundary summary preserves the claimed result in every admitted neighboring context, and separate any unbounded-time progress assumption from the safety induction.

\paragraph{First mathematical probe.} A bound on the number of live commitments is not a bound on the information needed to settle them. For the intended exact summary class, construct or rule out growing sets of pairwise contextually distinguishable boundary histories, keeping the local parameters fixed. If $N$ distinct histories must be distinguished by every such summary, it needs at least $N$ possible summary states (at least $\lceil\log_2 N\rceil$ bits in a fixed-width encoding). The required witness is a permitted continuation or neighboring context giving a different consumed decision, realizability verdict or preservation result; history length alone is not a distinction. Specify whether the summary is only a certificate interface or also a locally computable runtime account. Where a residual-profile quotient is used, justify its tests, adequacy under recurrence and comparison along actual source maps. Finite index alone does not establish effective computation, a strategy selector or complete synthesis. If the bounded-summary target fails, identify an informative positive subclass and the exact weakening (summary growth, information restriction or changed policy class) required. A lower bound plus a positive bounded subclass may be a substantial shared obstruction, but is not the full moonshot.

Attempt reduction to \href{https://doi.org/10.4230/LIPIcs.FSTTCS.2014.639}{acyclic rendezvous synthesis} and \href{https://lmcs.episciences.org/736}{parameterized token passing}; compare the exogenous two-process restriction with \href{https://arxiv.org/html/2611.04561v1\#S1.SS1}{Bollig et al.'s regular message-adversary theorem}. The latter is a frontier anchor for fixed legal delivery histories, not a result about action-dependent handoffs. Also check dynamic hierarchical-information synthesis and existing imperfect-information games when their actual architecture and policy assumptions apply. Geometric acyclicity is not an information-hierarchy theorem. Validate that a reduction neither grants new observation or authority nor eliminates the action-dependent handoff. Use L1--L3 only where they compose with a complete decision method. The core integration obligations do not wait for the moonshot.

\section{Research order and programme judgment}\label{sec:order}
Establish a common versioned semantic reference and inventory crosswalk first, before separate research tracks rely on shared meanings or assign coverage. This initial reference fixes the meanings against which particular proofs, comparisons, implementations and inventory dispositions are attributed; it does not prohibit revision. When those meanings change, identify the affected results and supply the compatibility interpretation or recheck needed to carry them to the revised version. Keep unresolved incompatibilities explicit. These are evidence and compatibility obligations, not a prescribed process or allocation of power for amending the contract. Identify the closest imported results against that reference. Develop a small executable common witness with overlapping live handoffs, local evidence, actual ownership and a changed receiver dependency; use it as a regression and an early joint test of R1, R2 and cooperative R5. This initial witness need not become the eventual benchmark or the only formal family. In parallel with that calibration, select principal I1 and I2 targets against the strongest prior mechanisms. For each, record the unresolved capability, the family and budgets, the matched comparator and the evidence that would count as a material advance before building a large benchmark. R2 is the initial candidate for the repair nucleus, not an assumption that every mechanism must be designed around it. Develop R3, R4 and R6 with the interfaces they actually consume and supply, including physical feasibility feedback to the continuation problem; revise the shared abstraction where those tests expose a lost distinction. Inventory work and moonshot obstruction analysis may proceed in parallel. Reproduce relevant known comparisons before assigning novelty claims. Build both integration returns from the same stack rather than two demonstrations.

Judge the results on separate axes. Scientific discharge requires both the coherent shared return in Section~\ref{sec:laws} and the substantive research obligations at their stated scope. A common new result may meet several targets through demonstrated transfer; inherited ingredients may support those proofs, but inherited coverage alone does not discharge an R1--R5 extension requirement. R6 retains its synthesis exception. Specialized advances remain valuable where stronger transfer fails, with the boundary and their consequential connection to the programme recorded. Engineering success requires I1 and I2, their positive service envelope, and actual constructive guarantee reuse, whether inherited or novel. Each R1--R6 target and its materially distinct subclaims receive an explicit result or unresolved disposition at their actual scope; partial and open obligations are not described as discharged. A full return establishes both the shared scientific and integrated engineering results, including the consequential connection of a new shared result to the stack, and resolves the substantive track targets. An otherwise valuable partial return must say which targets remain open. The moonshot has its own independent verdict. Before claiming closure, independent review should attack the weakest premise with a matched counterexample, replay the evidence and state what remains unresolved.

This is a developing research contract. It does not select implementations, fix numerical envelopes or tests, prescribe a process for changing the contract, or authorize treating a promising hypothesis as a result. The research questions and organizing laws may change as the mathematical and engineering boundaries become known. The evidentiary freeze of particular comparisons is distinct from governance of the research programme itself.

\appendix
\section{Sources and scope crosswalk}\label{sec:sources}
The related-project references in Section~\ref{sec:imports} are pinned to repository \href{https://github.com/eternal-november/eos-research-programmes/commit/4359a8a1484498b9c46d8e58a0112448a1fba06b}{commit \texttt{4359a8a14844}} (full revision in the links). Use the \repo{contextual-composition/NAVIGATION.md}{Contextual Composition navigation} and \repo{byzantine-control/README.md}{Byzantine Control index} to locate actual proof statements, executable sources and known limits; source claims retain their individual proof status and hypotheses. Inspect further or newer results when they become relevant. These routes are not an exhaustive novelty survey or a substitute for replaying a consumed proof. The references below retain the v0.1 prior-evidence and inventory crosswalk.
Use the \href{https://github.com/eternal-november/eos-research-programmes/blob/main/projects/swarm-coordination/research/COMPLETION.md}{v0.5 evidence map} to locate completed proofs and negative controls. The correspondence below locates prior material and names the initial law dependencies. It is not a partition or proof of transfer; each obligation has its own acceptance conditions.
\begin{description}
\item[R1:] L1--L2; C1a (communication) and C1b (local perception). Start from the \href{https://github.com/eternal-november/eos-research-programmes/blob/4359a8a1484498b9c46d8e58a0112448a1fba06b/projects/swarm-coordination/research/communication/Communication_Regimes.tex}{communication and observation notes}.
\item[R2:] L1--L3; C2 (rejoin), C3 (roles), C4 (partner substitution), C10 (live synthesis), and C11 (propagation and proof transfer).
\item[R3:] L1, L3; C5 (actions that create evidence) and C7 (forecasts and preparation).
\item[R4:] L2--L3; C6 (joint risk), C12 (goal revision), and C13 (risk revision).
\item[R5:] L1--L3; C8 (several owners and limited trust) and C9 (conflict and agreements); read alongside R2 for transfers and R4 for changes to authorized objectives.
\item[R6:] L1--L3; added formation track. Its exception applies only to a claim of innovation, not to evidence, safety, or service.
\item[I1:] L1--L3; added reference architecture. Use both the \href{https://github.com/eternal-november/eos-research-programmes/blob/4359a8a1484498b9c46d8e58a0112448a1fba06b/projects/swarm-coordination/fault-inventory/Swarm_Coordination_Fault_Inventory.tex}{mechanism inventory} and the \href{https://github.com/eternal-november/eos-research-programmes/blob/4359a8a1484498b9c46d8e58a0112448a1fba06b/projects/swarm-coordination/fault-inventory/Swarm_Fault_Inventory_Cleanroom.tex}{cleanroom inventory}. Use Byzantine Control for applicable laws, source interpretations, fault selection and regression witnesses, not as a substitute for either inventory.
\item[I2:] L1--L3; added independently selected enclosed-space integration. No single simulation can establish I1 inventory coverage.
\item[Moonshot:] L1--L3; targets uniform instance synthesis over corridor sizes and unbounded sequences of handoffs; it does not replace the finite-model engineering return.
\end{description}

\section{Conservative-extension audit}\label{sec:audit}
The audit below first records the refinements in v0.3.1 and then the additions in v0.3. The remaining audit is retained from v0.2.1 and compares the supplied v0.1 and v0.2 TeX sources with v0.2.1. The substantive baseline is v0.1; the shared laws, semantic distinctions, imported-result boundaries and integration safeguards introduced in v0.2 are retained. The permissive v0.2 clauses identified below are replaced. No research family, integration obligation, required experimental discriminator or original source link is removed. This is an editorial scope audit, not evidence that the proposed research has been completed or that the cited proofs have been replayed in this revision.

\subsection*{Changes in v0.3.1}
This patch retains the complete research and integration burden of v0.3, with the following refinements. The acceptance section separates established-mechanism advancement from comparison of alternative implementations of a new mechanism and from dependency ablations. It distinguishes operational capability, assurance and efficiency or service bounds; an assurance-only gain cannot substitute for the required positive integration advance. I1 makes source-specific adaptation, reused decision and proof work, and the actual cause of failed composition explicit. I2 requires a common consequential dependency for the physical and informational changes and tests the separate and joint contributions of coordination and recovery where admissible, with incompatible variants explicitly accounted for. Additional literature is supplied as reduction leads, not as established reductions. The research order restores freedom to replace or supplement the initial small witness. These are evidentiary and scope refinements or added discriminators, not a new research track, a fixed numerical metric or an amendment-governance procedure. R1--R6, L1--L3, the common semantic obligations, the mandatory Contextual Composition application and the moonshot are unchanged; all v0.3 source links and original integration cases are retained.

\subsection*{Changes in v0.3}
I1 adds a family-level constructive composition guarantee, a failure of independently correct recovery handlers, a frozen transfer that changes a recovery dependency or resource type, and a matched architectural advance. I2 adds a family-level coordination advance with a positive construction, a discriminating boundary, and comparisons to underground coordination and robust route execution. I2 still evaluates the I1 stack on held-out geometry. I1 inventory coverage, the original failure and transfer controls, R1--R6, the Contextual Composition obligation and the moonshot are retained. This change does not assert historical novelty, require physical tests, or fix a solver or numeric service floor.

\subsection*{Material revisions and their justification in v0.2.1}
The status labels distinguish restoration of an earlier obligation from an added acceptance condition or a clarification. They do not claim formal equivalence between the prose versions.

{\small
\renewcommand{\arraystretch}{1.15}
\begin{longtable}{@{}P{0.16\linewidth}P{0.13\linewidth}P{0.64\linewidth}@{}}
\toprule
Area & Status & Present obligation and reason \\
\midrule
\endfirsthead
\toprule
Area & Status & Present obligation and reason \\
\midrule
\endhead
Research closure & Restored & R1--R5 retain their research-extension burdens. Established coverage may supply engineering and redirect research, but does not automatically close an extension. One new result may serve several targets through genuine transfer. R6 alone retains the synthesis/comparison exception. This removes the weaker v0.2 inherited-track closure. \\
Comparison and negative results & Restored & Nearest-result analysis covers every obligation and materially distinct subclaim. Scoped lower bounds remain partial unless they resolve the target. A generic obstruction cannot replace a separately required construction, positive case, experiment or service guarantee. \\
Contextual Composition & Restored & At least one substantive application of a named preservation or replacement law is required. Duties and observations allowed across the interface, concrete hypotheses and a reused result are explicit. The direct comparison may tie or win; a categorical runtime is not required. \\
Multi-owner regimes & Restored and made explicit & Cooperative fleets lack a common controller. Limited trust remains distinct even with compliant owners and consistent truthful reports. Faulty/inconsistent evidence and competing objectives retain their separate burdens. The distinction prevents a centralized arbiter or Byzantine fault budget from silently replacing the original multi-owner problem. \\
Common semantics & Restored, with revision allowed & A versioned reference and inventory crosswalk precede reliance by separate tracks. Changed meanings need an identified compatibility interpretation or recheck. Permanent semantic immutability and amendment governance are not imposed. \\
Programme coherence & Strengthened & Every family has a consequential contribution, dependency or obstruction in the common body, with actual interpretations and compatible consumers. One isolated two-family success, a shared record or conjunction of unrelated results is insufficient. Specialized laws remain where transport fails; no universal theorem or theorem quota is imposed. \\
Moonshot mechanism & Restored & The synthesis procedure assembles boundary summaries into local controllers, and same-summary replacement preserves the result in admitted neighboring contexts. An unrelated whole-system solver plus unused summaries is insufficient. \\
Moonshot claim scope & Clarified & Uniform instance synthesis is distinct from one controller template for all sizes. A restricted policy class needs task justification; completeness within it does not prove impossibility for more expressive admissible controllers. A narrowed result retains its narrower scope. \\
Evaluation and amendment & Clarified & Evidence used for confirmation or closure retains the preselection freeze. Exploration and revision remain possible, but exploratory or post-change evidence cannot be relabelled as independent confirmation of the earlier claim. No rule determines who may amend the contract or how amendments are approved. \\
\bottomrule
\end{longtable}
}

\subsection*{Retained obligations and discriminators}
The following ledger locates the material v0.1 obligations in the present sections. Each entry retains its complete detailed requirements in the body; this ledger is not a replacement specification.

{\small
\renewcommand{\arraystretch}{1.15}
\begin{longtable}{@{}P{0.12\linewidth}P{0.81\linewidth}@{}}
\toprule
Scope & Preserved burden \\
\midrule
\endfirsthead
\toprule
Scope & Preserved burden \\
\midrule
\endhead
Governing scope & Safety and useful service under changing observations, capability, rights, roles and goals; prior v0.5 evidence and ordinary-method calibration; civil proofs, executable models and simulation; no physical tests, deployment claim or earlier private vehicle input. \\
Acceptance & Strongest applicable reductions; explicit inheritance and increments; actual local observations and powers; matched information, authority, resources and deadlines; charged preprocessing, retained state and runtime costs; safety separate from service; reproducible sources, counterexamples and appropriate proof kernels. See Sections~\ref{sec:semantics}, \ref{sec:imports} and \ref{sec:acceptance}. \\
R1 & Two evidence cuts into one degraded regime; one-way, delayed, partitioned and absent contact, including no usable observation or prior exchange; held-out non-radio substitution; anonymous/intermittent vision versus persistent identity or interlock; reception, action evidence and exclusion distinguished. \\
R2 & Overlapping handoffs, delayed recognition and live duties; changed later evidence or authority; roles, partners and split/rejoin without reset; dependency-sensitive propagation; affected-work/state bounds; frozen receiver with a held-out consumed dependency change; cached and whole-system comparisons; bounded synthesis with a valid late/failure safety path and charged service loss. \\
R3 & Action changes another actor's lawfully obtainable evidence; passive and strong active baselines; accurate but late/private evidence; a separate identification/forecast target with a non-identifiable pair and sharp bound or separating intervention; fixed-experiment reproduction; pre-fault change and false-forecast cost. \\
R4 & Scalable overlapping exposures; exact local admission or sharp sharing limit; one fixed uncertain mission law and the original any-failure cap; no fresh per-decision caps or independent worst cases; strong factored/common-information baselines; prospective and live-exposure amendments, timely/late notice, irreversible commitments and service feasibility. \\
R5 & Rights-preserving amendments change future amendment powers; decentralized cooperation, limited trust, inconsistent/faulty evidence and conflict; refusal, inspection, override, settlement and human takeover; fixed-consent baseline; changed future powers affect realizability; global agreement with locally impossible lawful execution. \\
R6 & Competing geometry, density, coverage, connectivity and clearance requirements; throughout invariants, deadline restoration and authorized relaxation; joint feasibility/recovery guarantee with a matched infeasible case; physical versus local-evidence infeasibility; old shape optional; synthesis exception without waiving safety or service. \\
I1 & Both inventories with original identifiers and all entry types; supported operational claims or validated exclusions/impossibilities with residual duties and lost service; open entries and assurance evidence; actual source mechanisms; second fault, common cause, masking and stale authorization; independent compound holdout; one stack, two materially different missions, reused safety and service, matched factored direct comparison. \\
I2 & Same stack; mandatory civil work, report deadlines and egress/fallback; separate report custody and vehicle recovery; known then learned branching geometry with routes/contact/sight/maps distinct; matched cave planner; frozen stack and independently selected geometry/order; blocked travel plus lost contact plus second fault; continued work and boundary refusal; resource and service metrics; trapping boundary and renewed holdout after case-specific changes. \\
Moonshot & Every positive corridor size; overlapping neighboring handoffs; bounded live commitments; action-dependent clearance authority; prescribed nonvacuous service; unbounded duration with no global reset or synchronized global observation; uniform laws and locally bounded summaries; memory/communication/deadline scaling; same-summary replacement; progress separate from safety; original reductions and separate verdict for a bounded subclass. \\
Order and sources & Common semantic reference and inventory crosswalk first; known comparisons before novelty claims; one integration stack; independent challenge and replay before closure; all original source links retained, with C1--C13 mappings in Section~\ref{sec:sources}. The moonshot does not gate core integration. \\
\bottomrule
\end{longtable}
}

\begin{samepage}
\paragraph{Retained v0.2 additions.} L1--L3 remain organizing hypotheses, with source-to-law interpretations and actual architectural consequences. The shared semantics retains local-policy quantifiers, maximal histories, actual effects, powers and source dependence. The track-specific additions retain dependency discovery and migration costs, decision-relevant identification, ex ante versus conditional risk, authority versus effective influence, and viability versus instantaneous feasibility. Integration retains a positive service envelope, source-history and model/implementation correspondence, separation of evaluator knowledge from controller inputs, consequential decision traces, and fair direct comparators distinct from ablations. The moonshot retains its distinguishability probe and the distinction between finite index, effective summaries and complete strategy synthesis. These additions strengthen the research contract; they are not presented as requirements already explicit in v0.1.

\paragraph{Open design choices.} The final law statements, implementations, numerical envelopes, test families and permitted configurations remain to be developed. The research contract's amendment governance remains unspecified; R5 still studies authority and amendment laws for the operational systems. Versioned evidence and honest dispositions constrain what a result can establish; they do not supply an approval process, a veto allocation or an immutable constitution for the programme.
\end{samepage}

\end{document}
